\chapter{Formalizzazione del Processo di Analisi Funzionale}
\section{Introduzione}
Il sistema di analisi è modellato come una sequenza di trasformazioni tra domini tipati. La pipeline complessiva è definita dalla composizione di tre fasi principali: Estrazione Sintattica, Risoluzione Semantica e Costruzione del Grafo. 
Siano $\mathcal{L}$ l'insieme dei linguaggi supportati, $\mathcal{S}_{code}$ il dominio del codice sorgente e $\mathcal{D}_{IR}$ il dominio della rappresentazione intermedia definito nel capitolo precedente.

La funzione di analisi completa $\Phi$ è definita come la composizione delle tre fasi fondamentali:
$$\Phi : \mathcal{L} \times \mathcal{S}_{code} \to \mathcal{G}$$
$$\Phi = G \circ R \circ E$$
Dove:
\begin{itemize}
    \item $E : \mathcal{L} \times \mathcal{S}_{code} \to \vec{\mathcal{M}}$ è la funzione di \textbf{Estrazione}.
    \item $R : \vec{\mathcal{M}} \to \vec{\mathcal{M}}$ è la funzione di \textbf{Risoluzione} dei riferimenti.
    \item $G : \vec{\mathcal{M}} \to \mathcal{G}$ è la funzione di \textbf{Costruzione del Grafo}.
\end{itemize}
Questa formulazione evidenzia come il sistema sia una pipeline pura dove l'output di ogni fase è l'input della successiva, garantendo la coerenza semantica attraverso il sistema dei tipi condiviso.

\section{Fase 1: Estrazione Sintattica ($E$)}
L'estrazione mappa i nodi dell'albero sintattico concreto ($\mathcal{N}$) in istanze del modello IR. Questa fase è l'unica dipendente dal linguaggio $L \in \mathcal{L}$. Poiché i nodi del CST contengono solo metadati strutturali, ogni funzione riceve in input anche il codice sorgente $S \in \mathcal{S}_{code}$ per l'estrazione dei token testuali.

\subsection{Funzioni Core}
\begin{itemize}
    \item \textbf{Generic Extract}: $E : \mathcal{L} \times \mathcal{S}_{code} \to \mathcal{M}$ \\
    Mappa l'intero sorgente in un unico \textbf{Modulo Radice} $\mathcal{M}$ che funge da entry-point del Component Tree.
    \item \textbf{Node Dispatcher}: $f_{disp} : \mathcal{N} \times \mathcal{S}_{code} \to \mathcal{C}^?$ \\
    Dato un nodo sintattico, identifica il costrutto e delega l'estrazione al parser specifico, restituendo un Componente opzionale.
\end{itemize}

\subsection{Parser Strutturali}
Le funzioni di mapping che istanziano i componenti principali:
\begin{itemize}
    \item $f_{\mathcal{S}} : \mathcal{N} \times \mathcal{S}_{code} \to \mathcal{S}^?$ (Tipi Strutturati)
    \item $f_{FN} : \mathcal{N} \times \mathcal{S}_{code} \to FN^?$ (Funzioni Libere)
    \item $f_{\mathcal{B}_{lk}} : \mathcal{N} \times \mathcal{S}_{code} \to \mathcal{B}_{lk}^?$ (Blocchi di Implementazione)
\end{itemize}

\subsection{Helper di Estrazione}
Funzioni atomiche che estraggono i tipi base necessari per la costruzione dei componenti dell'IR:
\begin{itemize}
    \item $\text{ext}_{\mathcal{QN}} : \mathcal{N} \times \mathcal{S}_{code} \to \mathcal{QN}^?$ \\
    Estrae il nome qualificato (es. \texttt{std::vec::Vec}) analizzando l'identificatore o il percorso sintattico.
    \item $\text{ext}_{\tau} : \mathcal{N} \times \mathcal{S}_{code} \to \tau$ \\
    Mappa un nodo tipo in un riferimento $\tau$ (inizialmente \texttt{Unknown} o \texttt{UserDefined}).
    \item $\text{ext}_{\vec{F}} : \mathcal{N} \times \mathcal{S}_{code} \to \vec{F}$ \\
    Estrae l'insieme dei campi di un tipo strutturato.
    \item $\text{ext}_{\vec{P}} : \mathcal{N} \times \mathcal{S}_{code} \to \vec{P}$ \\
    Estrae la lista dei parametri di una funzione, preservando ordine, tipi e variadicita.
    \item $\text{ext}_{\vec{\tau}_{sup}} : \mathcal{N} \times \mathcal{S}_{code} \to \vec{\tau}$ \\
    Estrae i riferimenti alle gerarchie di ereditarietà o implementazione di trait.
\end{itemize}

\section{Fase 2: Risoluzione dei Nomi e dei Tipi ($R$)}
La risoluzione trasforma un IR "parziale" (con riferimenti non risolti) in un IR semanticamente completo. Questa fase è totalmente \textbf{language-agnostic}.

\subsection{Tabelle dei Simboli e Contesto}
Sia $\Sigma = \mathcal{QN} \to \mathcal{C}$ la tabella dei simboli che associa nomi qualificati a componenti. Definiamo il contesto di risoluzione $\Gamma$ come la tripla $\langle \Sigma, \mathcal{QN}_{curr}, \vec{\mathcal{I}}_{imports} \rangle$.

\subsection{Funzioni di Risoluzione}
\begin{itemize}
    \item $\text{build\_sym\_table} : \vec{\mathcal{M}} \to \Sigma$ \\
    Scansiona l'albero dei componenti per popolare la tabella globale dei simboli.
    \item $\text{resolve}_{\tau} : \Gamma \times \tau \to \tau$ \\
    Data un riferimento $\tau$, cerca la corrispondenza univoca in $\Sigma$ applicando le regole di scoping definite in $\Gamma$.
\end{itemize}

\subsection{Ricorsione Strutturale}
La risoluzione è propagata ricorsivamente su tutta la gerarchia:
$$\text{resolve}_{\mathcal{M}} : \Gamma \times \mathcal{M} \to \mathcal{M}$$
Per ogni $\mathcal{S} \in \mathcal{M}$, viene invocata $\text{resolve}_{\mathcal{S}}$, che a sua volta invoca $\text{resolve}_{\tau}$ su ogni campo, parametro e super-tipo.

\section{Fase 3: Costruzione del Grafo delle Dipendenze ($G$)}
L'ultima fase mappa l'albero dei componenti risolti in un grafo orientato.

\begin{itemize}
    \item $\text{build\_graph} : \vec{\mathcal{M}} \to \mathcal{G}$
    \item $\text{extract\_edges} : \mathcal{C} \to \vec{E}$ \\
    Per ogni componente, genera archi $e = \langle \text{from}, \text{to}, k \rangle$ basandosi sui riferimenti $\tau$ contenuti (es. il tipo di un campo genera un arco \textit{UsesFieldType}).
\end{itemize}

\section{Proprietà del Modello}
La formalizzazione dimostra tre proprietà fondamentali del sistema:
\begin{enumerate}
    \item \textbf{Type Safety}: Ogni funzione $f: I \to O$ garantisce che l'output appartenga al dominio IR, impedendo stati inconsistenti tra le fasi.
    \item \textbf{Agnosticismo}: La separazione netta tra la Fase 1 ($L$-dependent) e le Fasi 2 e 3 ($L$-independent) assicura che il core logico non sia influenzato dalla sintassi del linguaggio sorgente.
    \item \textbf{Terminazione}: Poiché tutte le funzioni operano per ricorsione strutturale su alberi di sintassi finiti, il processo di analisi converge sempre in tempo finito.
\end{enumerate}

